\clearpage\SourceLeaf{207}
\begin{center}
[ON THE FIRST PRINCIPLES OF GEOMETRY.]
\end{center}

\begin{center}
\textsc{Gauss} to \textsc{Bessel}. Göttingen, 27 January 1829.
\end{center}

.\ .\ .\ .\ .\ .\ .\ From time to time, in a free hour, I have also thought again about another subject that has occupied me for nearly 40 years.
I mean the first principles of geometry. I do not know whether I have ever spoken to you
about my views on it. Here too I have been able to
consolidate some things further, and my conviction that we cannot ground geometry
completely a priori has grown stronger still, if that is possible.
It will probably be a long time before I get to prepare my \textls{very extensive} investigations of the subject for publication, and perhaps this will never happen in my lifetime,
for I dread the clamor of the Boeotians that would follow if I were to state my view in full. --- Strangely enough, \textls{besides} the familiar gap in
Euclid's geometry, which people have so far tried in vain to fill and will never
fill, there is another defect in it that, to my
knowledge, no one has yet pointed out, and that is by no means easy
(though possible) to remedy. This is the definition of the \textls{plane} as a surface
in which the straight line joining \textls{any two} points lies \textls{entirely}. This
definition contains \textls{more} than is needed to determine the surface and \emph{tacitly} involves a \textls{theorem} that must first be proved. .\ .\ .\ .\ .\ .

\rule{4cm}{0.5pt}

